


\chapter{光的偏振}

\begin{definition}[普通光源发出自然光，单一光波列是偏振光]
普通光源发光不是一个连续、稳定、方向固定的过程，而是大量原子或分子在能级之间跃迁时发出的。原子从高能级跃迁到低能级时，会放出一段电磁波。这一小段长度有限的光波叫一个波列。单个原子一次跃迁发出的一段波列，它的电场振动方向在这段波列中可以是确定的，因此单个波列可以看作偏振光。但普通光源包含无数这样的波列，它们的方向、相位都乱，所以整体表现为自然光。
\end{definition}

\begin{definition}[偏振（Polarization）]
    光是电磁波。我们通常说光的振动方向，主要指电场矢量 \(\vec E\) 的振动方向。光沿某方向传播，比如沿 \(z\) 方向传播，那么电场 \(\vec E\) 一定在垂直于传播方向的平面内振动，也就是在 \(xOy\) 平面内振动。如果电场振动方向在这个横向平面中各方向都一样，没有哪个方向占优势，那就是自然光。如果电场振动偏向某个方向，或者只沿某个方向振动，那就是偏振光。只有横波才有偏振现象，纵波不发生偏振。
    \begin{itemize}
    \item 完全偏振光（线偏振光）：光振动只沿某一固定方向。传播方向和电场振动方向共同确定的平面，就是振动面。
    \item 部分偏振光：部分偏振光介于自然光和线偏振光之间。各方向都有，但某些方向更强，某些方向更弱，因此有方向偏好。
\[
\boxed{\text{部分偏振光}=\text{自然光成分}+\text{线偏振光成分}}
\]
    \end{itemize}
\end{definition}

\begin{definition}[偏振片（起偏器）]
    偏振片定义为：涂有二向色性材料的透明薄片。它具有二向色性，强烈吸收某一方向的电场振动，只允许另一个方向（这个方向叫此偏振片的偏振化方向）的电场振动通过。自然光可以等效地分解成两个相互垂直、互相独立、没有确定相位关系的等强度线偏振光。每个方向平均分得一半强度：
\[
I_x=I_y=\frac{I_0}{2}
\]偏振片只让沿偏振化方向的那个分量通过，于是透射光强变为：
\[
I=\frac{I_0}{2}
\]同时，透射光变成线偏振光。所以我们可以通过偏振片区分自然光、线偏振光和部分偏振光：
\[
\begin{array}{c|c}
\text{旋转偏振片后的现象} & \text{入射光类型}\\
\hline
I\text{不变} & \text{自然光}\\
I\text{变且有消光} & \text{线偏振光}\\
I\text{变但无消光} & \text{部分偏振光}
\end{array}
\]
\end{definition}

\begin{theorem}
自然光先通过起偏器，变成线偏振光。然后这束线偏振光再通过检偏器。设进入检偏器前的线偏振光强为 \(I_0\)，电场振幅为 \(E_0\)，检偏器的偏振化方向与入射线偏振光的振动方向夹角为\(\alpha\)，那么电场只有沿检偏器偏振化方向的分量能通过。所以透过检偏器后的振幅为：
\[
E=E_0\cos\alpha
\]因为光强正比于振幅平方\(I\propto E^2\)于是：
\[
\frac{I}{I_0}=\frac{E^2}{E_0^2}\implies I=I_0\cos^2\alpha
\]注意这里的 \(I_0\) 是入射到检偏器上的线偏振光强度，不是最初自然光的强度。如果最初自然光强为 \(I_s\)，先通过第一个偏振片后变为\(I_0=\frac{I_s}{2}\)，再通过第二个偏振片：
\[
I=I_0\cos^2\theta=\frac{I_s}{2}\cos^2\theta
\]
\end{theorem}

\begin{exercise}[两片偏振片]
两个偏振片叠加放在一起，强度为 \(I_0\) 的自然光垂直入射其上，若通过两个偏振片后的光强为 \(I_0/8\)，则此两偏振片的偏振化方向间的夹角（取锐角）是？若在两片之间再插入一片偏振片，其偏振化方向与前后两片的偏振化方向的夹角（取锐角）相等，则通过三个偏振片后的透射光强度为？
\end{exercise}
\begin{solution}
\[\frac{I_0}{8}=\frac{I_0}{2}\cos^2\alpha\implies \cos\alpha=\frac12\implies \alpha=60^\circ\]
原来两片相差 \(60^\circ\)，中间一片放在角平分方向，所以每相邻两片夹角为 \(30^\circ\)。
\[I_3=I_0\cdot\frac12\cdot\cos^2 30^\circ\cdot\cos^2 30^\circ=I_0\cdot\frac12\cdot\frac34\cdot\frac34=\frac{9I_0}{32}\]
\end{solution}


\begin{exercise}[2010级期末：三片偏振片中间片成 \(30^\circ\)]
三个偏振片 \(P_1,P_2,P_3\) 堆叠在一起，\(P_1\) 与 \(P_3\) 的偏振化方向互相垂直，\(P_2\) 与 \(P_1\) 的偏振化方向夹角为 \(30^\circ\)。强度为 \(I_0\) 的自然光垂直入射到 \(P_1\)，并依次透过 \(P_1,P_2,P_3\)。求最终透射光强。
\end{exercise}
\begin{solution}
\[\frac12I_0\cos^230^\circ\cos^260^\circ=\frac{3}{32}I_0\]
\end{solution}
